%----------------------------------------------------------------------------------------
% Tailored CV — Data Manager*in, ZPM Heidelberg
% Pathologisches Institut, Zentrum für Personalisierte Medizin
% Job-ID V000015540
%----------------------------------------------------------------------------------------

\documentclass[10pt,a4paper,sans]{moderncv}

\usepackage[utf8]{inputenc}
\usepackage[T1]{fontenc}

\moderncvstyle{classic}
\moderncvcolor{blue}

\usepackage[scale=0.78]{geometry}
\usepackage{ragged2e}

%----------------------------------------------------------------------------------------

\firstname{Kundai Farai}
\familyname{Sachikonye}

\address{Auf der Lichtung 19}{82178 Puchheim, Germany}
\mobile{(+49) 1626386344}
\email{kundai.sachikonye@bitspark.com}
\homepage{github.com/fullscreen-triangle}{github.com/fullscreen-triangle}

%----------------------------------------------------------------------------------------

\begin{document}

\makecvtitle

%----------------------------------------------------------------------------------------
%	EDUCATION
%----------------------------------------------------------------------------------------

\section{Education}

\cventry{2019--2021}{PhD candidate, Computational Lipidomics}{Technische Universität München / Bayerisches Forschungsinstitut}{}{}{
  ML for large-scale clinical omics cohorts; multi-omics data integration
  across federated clinical sites; distributed pipeline maintenance.
  Departed to pursue independent research.
}
\cventry{2018--2019}{MSc Computational Biology}{Universität Tübingen}{}{}{
  Probabilistic graphical models, statistical bioinformatics, Bayesian methods,
  computational systems biology, ML for biological sequences.
}
\cventry{2014--2017}{MSc Pharmaceutical Biotechnology}{HAW Hamburg}{}{}{}
\cventry{2011--2014}{BSc Biochemistry and Cell Biology}{Jacobs University Bremen}{}{}{}

%----------------------------------------------------------------------------------------
%	ONCOLOGY DATA INFRASTRUCTURE
%----------------------------------------------------------------------------------------

\section{Oncology Data Infrastructure}

\cvitem{syndrome}{
  \textbf{Cancer Patient and Tumour Characterisation Framework.}
  Structured disease state representation $\mathbf{D}\in[0,1]^8$ capturing
  immune activation ($D_\mathrm{Imm}$), barrier integrity ($D_\mathrm{Bar}$),
  metabolic ($D_\mathrm{ATP}$), circadian ($D_\mathrm{Circ}$), and four
  additional dimensions; Bayesian $P(R\mid Q)$ disease trajectory inference;
  sparse $\ell_1$ LP for therapeutic target selection; AUC$=1.00$
  across 25/25 validation experiments.
  Neoantigen identification, MHC binding prediction (AUC\,$0.81$ on 2,847
  IEDB peptides), immune escape scoring.
  Browser-deployable, no server infrastructure required.
  Live: \href{https://syndrome-seven.vercel.app/tools}{syndrome-seven.vercel.app/tools}.
}

\cvitem{gospel}{
  \textbf{Clinical Genomics Data Pipeline.}
  $10^6$ variants/second; germline and somatic variant calling; RNA-seq
  quantification and differential expression; Bayesian network multi-omics
  integration (host genotype $\to$ metagenomics $\to$ transcriptomics);
  ClinVar, KEGG, Reactome API integration for variant annotation.
  Validated: gnomAD~v3.1.2, 1000~Genomes Phase~3, UK~Biobank.
  Ray distributed, Docker deployable.
  \href{https://github.com/fullscreen-triangle/gospel}{github.com/fullscreen-triangle/gospel}
}

\cvitem{lavoisier}{
  \textbf{Clinical Mass Spectrometry Data Pipeline.}
  Dual pipeline: numerical MS1/MS2 processing + CNN spectral classification;
  98.9\,\% NIST accuracy; $>$100\,GB mzML at clinical cohort scale;
  memory-mapped I/O for low-memory environments; HDF5 storage.
  Applicable to translational proteomics and metabolomics data management.
  \href{https://github.com/fullscreen-triangle/lavoisier}{github.com/fullscreen-triangle/lavoisier}
}

\cvitem{four-sided-triangle}{
  \textbf{HPC Inference and Data Processing Pipeline.}
  Rust core (15--50$\times$ speedup via PyO3); Ray-distributed parallel
  execution; FastAPI orchestration; 8-stage metacognitive pipeline with
  automated benchmarking. Applicable to ZPM batch data processing workflows.
  \href{https://github.com/fullscreen-triangle/four-sided-triangle}{github.com/fullscreen-triangle/four-sided-triangle}
}

%----------------------------------------------------------------------------------------
%	EXPERIENCE
%----------------------------------------------------------------------------------------

\section{Experience}

\cventry{2021--present}{Independent Computational Biology Researcher}{Bitspark GmbH}{Munich}{}{
  Design, implementation, and validation of clinical data infrastructure for
  oncology, multi-omics integration, and pharmacology. All systems publicly
  deployed with independently verifiable validation against established standards.
}

\cventry{2019--2021}{Computational Lipidomics}{TUM / Bayerisches Forschungsinstitut}{Munich}{}{
  Multi-omics data integration across federated clinical sites; distributed
  pipeline maintenance; ML for clinical MS cohort data. Python, Java.
}

\cventry{2017--2018}{Glycomics and Proteomics}{Universitätsklinikum Hamburg-Eppendorf}{}{}{
  Automated LC-MS/MS and MALDI-TOF pipeline operation; clinical sample
  processing and quality control; data annotation and documentation for
  clinical cohort studies; multi-centre sample logistics coordination.
}

%----------------------------------------------------------------------------------------
%	TECHNICAL SKILLS
%----------------------------------------------------------------------------------------

\section{Technical Skills}

\cvitem{Data Analytics}{
  R (statistical analysis, Bioconductor pipelines, ggplot2, tidyverse);
  Python (pandas, numpy, scipy, scikit-learn, PyTorch);
  statistical modelling, Bayesian inference, dimensionality reduction (PCA, UMAP);
  clinical cohort data quality assurance
}

\cvitem{Clinical Data}{
  Genomic data standards: VCF, BAM, FASTQ; gnomAD, ClinVar, KEGG, Reactome;
  multi-omics integration (genomics, transcriptomics, proteomics, metabolomics);
  clinical trial data workflows; GDPR-compliant multi-site data handling;
  patient data pseudonymisation; AnnData / HDF5 structured clinical data
}

\cvitem{Data Infrastructure}{
  Relational and document databases; REST API integration;
  Ray distributed processing; Docker and Kubernetes;
  memory-mapped I/O for large dataset handling;
  CI/CD pipeline maintenance; Git version control
}

\cvitem{Communication}{
  Technical documentation for clinical and non-technical audiences;
  scientific manuscript writing (formal mathematical apparatus,
  quantitative validation); interdisciplinary project coordination
  (clinicians, computational researchers, IT teams);
  user training and support for analytical tools
}

%----------------------------------------------------------------------------------------
%	LANGUAGES
%----------------------------------------------------------------------------------------

\section{Languages}

\cvitemwithcomment{English}{Native / Fluent}{}
\cvitemwithcomment{German}{Conversational}{Living in Germany continuously since 2011}

%----------------------------------------------------------------------------------------

\renewcommand{\listitemsymbol}{-~}

\end{document}
